\section{An arithmetic transfer: one torsion point for many analytic constraints}
\label{sec:torsion-transfer}

The most portable arithmetic ingredient we found is not the final RH
statement. It is a selection theorem buried in
\texttt{New/Congruences/Selmer/RootsOfUnity.lean}. Its root-counting argument
controls the divisibility of many polynomials at the same tuple of roots of
unity, with an order bound independent of their degrees. This section
extracts a slightly stronger budget from that argument and proves its
extension to bounded power series and infinite families. These are derived
results of this analysis; no claim of literature priority is made.

Fix a prime $p$. Write $v$ for the valuation on $\C_p$, normalized by
$v(p)=1$, and set $v(0)=+\infty$. Its ring of integers is
$\mathcal O_{\C_p}=\{z:v(z)\ge0\}$. For a finite real number $t$, define
$(a-t)_+=\max(a-t,0)$, with $(+\infty-t)_+=+\infty$.
For nonnegative numbers indexed by an arbitrary set, a sum means the
supremum of its finite partial sums. In particular, no countability
assumption is implicit in the statements below.

\subsection{The finite polynomial mechanism and its quantitative budget}

The source theorem
\texttt{New.Selmer.exists\_rootsOfUnity\_ellVal\_eval\_le\_add}
selects one tuple for a finite family of integral polynomials. Its
intermediate lemmas retain more information than its final statement.
We give the argument, including that information, so the analytic
extension has a precise mathematical starting point.

For an integral polynomial $g(T)=\sum_k a_kT^k$ and $r>0$, put
\[
 v_r(g)=\min_k\{v(a_k)+kr\}.
\]
This is the Gauss valuation of radius $p^{-r}$. It is additive on
products, and $v(g(y))\ge v_r(g)$ when $v(y)\ge r$.
Set
\[
 r_h=\frac1{\varphi(p^h)}\qquad(h\ge1).
\]
Every primitive $p^h$-th root of unity satisfies $v(\zeta-1)=r_h$.

\begin{lemma}[One coordinate, with an explicit choice of orders]
\label{lem:one-coordinate-torsion}
Let $\rho>0$, $0<\delta<1$, and $C\ge0$. Suppose that $H$ satisfies
\[
 C<\delta\sum_{\substack{1\le h\le H\\r_h\le\rho}}\frac1{h+1}.
\]
If $g\in\overline{\Z}_p[T]$ and $v_\rho(g)\le C$, there are an
admissible $h\le H$ and a primitive $p^h$-th root of unity $\zeta$ such
that
\[
 v(g(\zeta-1))\le v_{r_h}(g)+\delta.
\]
Here $\overline{\Z}_p$ denotes the integral elements of
$\overline{\Q}_p$.
\end{lemma}

\begin{proof}
For each admissible $h$, let $R_h$ be the multiset of roots of $g$ having
valuation $r_h$. Factorization over $\overline{\Q}_p$ and multiplicativity
of the Gauss valuation give
\[
 \sum_{\substack{1\le h\le H\\r_h\le\rho}}\#R_h\,r_h
 \le v_\rho(g)\le C.
\]
Indeed, a root $\alpha$ of positive valuation contributes
$\min(v(\alpha),\rho)$ to $v_\rho(g)-v_0(g)$, and it belongs to at most
one of the $R_h$. The remaining term $v_0(g)$ is nonnegative.
Consequently some admissible $h$ has
$\#R_h\,r_h\le\delta/(h+1)$.

If $\alpha\in R_h$, then $1+\alpha$ is a unit. For this $\alpha$,
\[
 \sum_{\zeta\text{ primitive of order }p^h}
      \min\{v(\zeta-1-\alpha),1\}\le h+1.
\]
To see this, use the polynomial
$q(Y)=(Y+1+\alpha)^{p^h}-1$. Its coefficient of $Y$ has valuation $h$,
so $v_1(q)\le h+1$. The numbers $\zeta-1-\alpha$ in the displayed sum
are among its roots; the root-factorization formula for $v_1(q)$ gives
the inequality. Sum over $R_h$ and average over the
$\varphi(p^h)=1/r_h$ primitive roots. Some $\zeta$ then satisfies
\[
 \sum_{\alpha\in R_h}\min\{v(\zeta-1-\alpha),1\}
 \le\#R_h\,r_h(h+1)\le\delta<1.
\]
No summand can equal $1$, so the minima can be removed and none of the
differences vanishes. Roots at other valuations contribute exactly their
Gauss-valuation terms. Hence
\[
 v(g(\zeta-1))
 =v_{r_h}(g)+\sum_{\alpha\in R_h}
       \bigl(v(\zeta-1-\alpha)-r_h\bigr)
 \le v_{r_h}(g)+\delta.
\]
This is the source's root-counting and averaging argument.
\end{proof}

Fix an upper bound $s_0\ge1$ on the number of variables, a positive
rational $\rho$, a total valuation budget $B\ge0$, and an allowed loss
$\epsilon>0$. Define
\begin{equation}
\label{eq:torsion-parameters}
 h_* = \min\{h\ge1:r_h\le\rho\},\qquad
 \delta=\frac{\min(\epsilon,1)}{s_0+2}.
\end{equation}
Choose $H$ so that
\begin{equation}
\label{eq:harmonic-budget}
 B+1<\delta\sum_{h=h_*}^{H}\frac1{h+1}.
\end{equation}
A completely explicit sufficient choice is
\begin{equation}
\label{eq:explicit-torsion-height}
 H=\left\lceil(h_*+1)
             \exp\!\left(\frac{B+1}{\delta}\right)\right\rceil.
\end{equation}
The integral comparison
$\sum_{h=h_*}^{H}(h+1)^{-1}\ge
\log((H+2)/(h_*+1))$ proves the strict inequality in
\eqref{eq:harmonic-budget}. This is a sufficient bound extracted from
the proof, not an asserted optimal bound. The maximum order $p^H$ can
be roughly double exponential in $1/\epsilon$; the result is primarily
a structural existence theorem.

\begin{proposition}[Finite polynomial selection with positive-excess control]
\label{prop:polynomial-excess}
With these parameters and $1\le s\le s_0$, let $P_i$ be a finite
family in $\overline{\Z}_p[T_1,\ldots,T_s]$. Suppose that algebraic
witness points $x_i$ satisfy
\[
 v(x_{ij})\ge\rho,\qquad
 t_i=v(P_i(x_i))<\infty,\qquad\sum_i t_i\le B.
\]
There is one tuple $\zeta=(\zeta_1,\ldots,\zeta_s)$ such that
\[
 \zeta_j^{p^H}=1,\qquad r_H\le v(\zeta_j-1)\le\rho,
 \qquad
 \sum_i\bigl(v(P_i(\zeta-1))-t_i\bigr)_+\le\epsilon.
\]
In particular all the selected values are nonzero.
\end{proposition}

\begin{proof}
Replace the coordinates one at a time, using the same replacement for
all polynomials. At a given step, free the coordinate being replaced
in each polynomial and freeze its other coordinates at its own current
point. This gives integral one-variable polynomials $g_i$. If $c_i$
are the current value valuations, their sum is at most $B+1$: after
$k\le s_0$ steps the total increase will be at most $k\delta\le1$.
The coordinate about to be replaced still has valuation at least
$\rho$, so
\[
 v_\rho\left(\prod_i g_i\right)
 =\sum_i v_\rho(g_i)\le\sum_i c_i\le B+1.
\]
Apply Lemma~\ref{lem:one-coordinate-torsion} to the product and let
$a_i=v(g_i(\zeta-1))$, $b_i=v_{r_h}(g_i)$. Its conclusion is
\[
 a_i\ge b_i,\qquad c_i\ge b_i,\qquad
 \sum_i a_i\le\sum_i b_i+\delta.
\]
The inequality $c_i\ge b_i$ uses $r_h\le\rho$. All these quantities
are finite. Therefore
\[
 \sum_i(a_i-c_i)_+
 \le\sum_i(a_i-b_i)\le\delta.
\]
In particular the total valuation grows by at most $\delta$, justifying
the inductive bound used above. The positive part of a sum of changes
is at most the sum of their positive parts. After all $s$ coordinates,
the total positive excess is at most $s\delta\le\epsilon$.

Every selected coordinate is a primitive root of an order $p^h$ with
$h_*\le h\le H$. This gives both stated radius bounds and the common
order bound. The estimates in the source's
\texttt{exists\_ellVal\_eval\_update\_le\_add} prove precisely the
three inequalities above; its final wrapper keeps the individual
conclusion and discards the stronger aggregate information.
\end{proof}

\subsection{Completion and infinitely many simultaneous constraints}

\begin{theorem}[Uniform torsion selection for bounded power series]
\label{thm:analytic-torsion-selector}
Keep the parameters in \eqref{eq:torsion-parameters}--
\eqref{eq:harmonic-budget}. Let $I$ be any index set, and let
\[
 F_i\in\mathcal O_{\C_p}[[T_1,\ldots,T_s]],\qquad
 x_i\in\C_p^s,\qquad 1\le s\le s_0.
\]
Assume
\[
 v(x_{ij})\ge\rho,\qquad
 t_i=v(F_i(x_i))<\infty,\qquad
 \sup_{J\subset I,\,J\text{ finite}}\sum_{i\in J}t_i\le B.
\]
Then a single tuple of roots of unity satisfies
\[
 \zeta_j^{p^H}=1,\qquad r_H\le v(\zeta_j-1)\le\rho,
 \qquad
 \boxed{\sup_{J\subset I,\,J\text{ finite}}
 \sum_{i\in J}\bigl(v(F_i(\zeta-1))-t_i\bigr)_+\le\epsilon.}
\]
Consequently every $F_i(\zeta-1)$ is nonzero and
\[
 v(F_i(\zeta-1))\le t_i+\epsilon,\qquad
 \sup_{J\text{ finite}}\sum_{i\in J}v(F_i(\zeta-1))\le B+\epsilon.
\]
The bound $H$ is independent of the cardinality of $I$, the complexity
of the series, and the coefficient fields of the series and witnesses.
\end{theorem}

\begin{proof}
Let $G_H$ be the finite set of tuples satisfying the two order and radius
conditions. It is nonempty because $H\ge h_*$. A coordinate has
$p^H-p^{h_*-1}$ possible values, so
\[
 |G_H|=(p^H-p^{h_*-1})^s.
\]

First fix a finite $J\subset I$. Choose
$M>B+\epsilon+\rho+1$ and a nonnegative integer $N$ with
$(N+1)r_H>M$. Truncate each $F_i$, $i\in J$, at total degree $N$,
obtaining $P_i$. At each witness and every point of $G_H-1$, all
coordinates have valuation at least $r_H$. Thus every omitted
monomial has valuation at least $(N+1)r_H$. There are only finitely
many monomials of each total degree; their valuations tend uniformly
to infinity as the degree grows. Completeness of $\C_p$ and its
nonarchimedean inequality therefore show that the series converge
there and
\begin{equation}
\label{eq:uniform-torsion-tail}
 v(F_i(y)-P_i(y))\ge(N+1)r_H>M
\end{equation}
at every relevant point $y$.

There are now finitely many polynomial coefficients and witness
coordinates. By density of $\overline{\Q}_p$ in $\C_p$, approximate
every coefficient of the $P_i$ by an integral algebraic number, with
error valuation greater than $M$, and call the resulting polynomials
$Q_i$. Approximate each witness coordinate similarly by $x'_{ij}$.
As $M>\rho$, the new witnesses still satisfy $v(x'_{ij})\ge\rho$.

Integral polynomials are $1$-Lipschitz for the maximum norm on the unit
polydisc. Indeed, telescope the difference of two products of integral
coordinates and then sum the monomials. Thus replacing these witness
coordinates changes polynomial values by valuation greater than $M$.
Replacing coefficients has the same precision everywhere on the
unit polydisc. Together with \eqref{eq:uniform-torsion-tail} this gives
\[
 v(Q_i(x'_i)-F_i(x_i))>M,
 \qquad v(Q_i(x'_i))=t_i,
\]
because $t_i\le B<M$.

Apply Proposition~\ref{prop:polynomial-excess} to the $Q_i$, $i\in J$.
It gives $\zeta\in G_H$ with positive-excess budget at most $\epsilon$.
Each selected polynomial value has valuation at most
$t_i+\epsilon\le B+\epsilon<M$. Its difference from
$F_i(\zeta-1)$ has valuation greater than $M$, by coefficient
approximation and the uniform tail bound. The ultrametric inequality
then gives the exact equality
\[
 v(F_i(\zeta-1))=v(Q_i(\zeta-1)).
\]
The desired budget follows for this finite $J$, without consuming any
additional error allowance.

For an arbitrary $I$, let $S_J\subset G_H$ consist of the tuples
satisfying the excess bound for the finite set $J$. Every $S_J$ is
nonempty. For finitely many sets $J_1,\ldots,J_k$,
\[
 S_{J_1\cup\cdots\cup J_k}\subseteq
 S_{J_1}\cap\cdots\cap S_{J_k},
\]
since all excess summands are nonnegative. The family $(S_J)$ has the
finite intersection property inside a finite set; its total
intersection is therefore nonempty. A tuple in that intersection
proves the theorem. Each singleton constraint rules out an infinite
valuation, and summing $v(F_i(\zeta-1))\le
t_i+(v(F_i(\zeta-1))-t_i)_+$ gives the final aggregate bound.
\end{proof}

\begin{remark}[What the completion step adds]
Over one fixed finite extension $K/\Q_p$, with ramification index $e$,
a series in $\mathcal O_K[[T]]$ whose constant term is not a unit has
valuation at every witness of depth $\rho$ at least
$\min(\rho,1/e)$. The total budget then permits only finitely many
such series; a unit constant term is nonzero throughout the open unit
polydisc. Thus allowing $\mathcal O_{\C_p}$ coefficients, or algebraic
coefficient fields with unbounded ramification, makes the
infinite-family conclusion materially stronger. The finite-family
analytic result already applies over a fixed $K$.
\end{remark}

\subsection{Matrices, cokernels, and the precise arithmetic application}

\begin{corollary}[Simultaneous inverse and denominator control]
\label{cor:torsion-matrices}
Let $A_i(T)$ be square matrices, of possibly different finite sizes,
with entries in $\mathcal O_{\C_p}[[T_1,\ldots,T_s]]$. Suppose the
witness values
$t_i=v(\det A_i(x_i))$ obey the hypotheses of
Theorem~\ref{thm:analytic-torsion-selector}. Then a single tuple in
$G_H$ makes all the matrices invertible and satisfies
\[
 \|A_i(\zeta-1)^{-1}\|_{\max}\le p^{t_i+\epsilon}.
\]
The same positive-excess and total budgets hold for their determinant
valuations.
\end{corollary}

\begin{proof}
Apply the theorem to $F_i=\det A_i$. Adjugate entries are integral at
the selected point, and
$A_i^{-1}=\operatorname{adj}(A_i)/\det A_i$.
\end{proof}

If the coefficients belong to a fixed $\mathcal O_K$, the evaluations
belong to the finite extension $L=K(\mu_{p^H})$. Smith normal form
over the discrete valuation ring $\mathcal O_L$ gives the exact identity
\[
 v(\pi_L)\,\operatorname{length}_{\mathcal O_L}
      \operatorname{coker}A_i(\zeta-1)
 =v(\det A_i(\zeta-1)).
\]
Thus the theorem controls normalized lengths, not merely
nonvanishing. For a rectangular presentation with $r$ rows, choose an
$r\times r$ minor that is nonzero at the witness and apply the theorem
to that minor. Its cokernel surjects onto the cokernel of the full
presentation, so its determinant valuation bounds the latter's
normalized length above.

This supplies a concrete transfer problem: construct integral analytic
presentation matrices for a family of Selmer modules or
projection-denominator modules, and prove the nonzero-witness budget.
Once those arithmetic inputs are supplied, one finite-order character
specialization can control many modules at once. The theorem does not
construct such a presentation or establish the budget for an arbitrary
elliptic curve.

\subsection{Boundaries that prevent false conjecture proofs}

\paragraph{Arbitrary analytic functions are too general.}
The function
\[
 \log(1+T)=\sum_{n\ge1}(-1)^{n+1}\frac{T^n}{n}
\]
converges for $v(T)>0$. It has a nonzero value at $T=p^2$, of
valuation $2$, but vanishes at every $T=\zeta-1$ with $\zeta$ of
$p$-power order. Its coefficient valuations $-v_p(n)$ are unbounded
below. No fixed denominator makes this an integral bounded series.
The coefficient hypothesis is therefore essential.

\paragraph{Infinitely many separate bounds do not give a total budget.}
For $h\ge1$, take
\[
 F_h(T)=\Phi_{p^h}(1+T)\in\Z_p[T],\qquad x_h=p^2.
\]
The identity
$\Phi_{p^h}(u)=(u^{p^h}-1)/(u^{p^{h-1}}-1)$ and the elementary
lifting-the-exponent calculation for $u=1+p^2$ show that
$v(F_h(p^2))=1$ for every $h$, including $p=2$.
Nevertheless each nontrivial $p$-power root annihilates one $F_h$.
A separate bound $t_h\le1$ cannot replace the finite total budget.

\paragraph{The order may vary inside the finite grid.}
A single polynomial $\Phi_{p^h}(1+T)$ rules out every root of exact
order $p^h$, while having a finite witness valuation. Thus the result
does not allow an arbitrary prescribed exact order or conductor.
Likewise the product over $1\le h\le H$ vanishes at every nontrivial
$p$-power root of order at most $p^H$ and has witness valuation $H$.
Dependence on the valuation budget cannot be discarded. No density
statement about successful characters is proved.

\paragraph{Small normalized length does not automatically mean zero.}
The character and its coefficient field change as $\epsilon$ changes.
In a ramified field the smallest positive normalized length is
$v(\pi_L)=1/e_L$, and $e_L$ may grow rapidly with $H$. Consequently
producing, for every $\epsilon>0$, a different specialization of
normalized length at most $\epsilon$ does not by itself produce a
zero Selmer module. One must also control the coefficient field or
prove a suitable uniform discreteness statement. This prevents a
spurious deduction of Selmer vanishing from arbitrarily small losses.

\subsection{A structural obstruction for the unmodified elliptic Tate module}

The source's general twisting theorem also has a restriction that
cannot be ignored. The structure \texttt{New.Selmer.SelmerDatum} in
\texttt{SelmerData.lean} requires an inertia element $\gamma_*$ at a
place away from $p$ that acts by a scalar $\eta$, with
$\eta^2-1$ a unit. This hypothesis is incompatible with an unmodified
rank-two elliptic Tate module.

\begin{proposition}[Scalar-inertia obstruction]
\label{prop:scalar-inertia-obstruction}
Let $\rho$ be a rank-two representation over a nonzero integral
domain. If $\det\rho$ is trivial on inertia away from $p$, then no
element of that inertia can act as $\eta I_2$ with $\eta^2-1$ a unit.
\end{proposition}

\begin{proof}
Taking determinants gives $1=\det(\eta I_2)=\eta^2$.
Hence $\eta^2-1=0$, which is not a unit.
\end{proof}

For an elliptic curve, the Weil pairing identifies
$\bigwedge^2T_pE$ with $\Z_p(1)$. The determinant is therefore the
cyclotomic character, which is unramified at places of residue
characteristic different from $p$. The proposition applies, as it
also does to its Tate twists and dual. A preliminary finite-order
character twist can change the determinant; relating the resulting
local conditions and Selmer groups back to the desired arithmetic
problem is additional work.

The most promising connection is therefore to determinant and
specialization methods in $p$-adic families, including the missing
cases of refined Kurihara and Iwasawa statements. Existing work already
proves simultaneous nonvanishing for many finite families. The
specific added tool here is a uniform finite grid and a controlled
total valuation loss for bounded analytic data. It does not establish
an Iwasawa determinant identity, create a nonzero Euler-system class
from no initial witness, remove a functional-equation forced zero, or
prove a positive proportion of nonvanishing twists.
